\documentclass[11pt]{article}

\usepackage[T1]{fontenc}
\usepackage{lmodern}
\usepackage{amsmath,amssymb,amsthm}
\usepackage[margin=1in]{geometry}
\usepackage{booktabs}
\usepackage{enumitem}
\usepackage{microtype}
\usepackage[hidelinks]{hyperref}

\hypersetup{
  pdftitle={General-m minimal-equilibrium stability: proved regions and a subcritical steady obstruction},
  pdfauthor={ShenWork formalization project}
}

\newcommand{\R}{\mathbb{R}}
\newcommand{\norm}[1]{\left\lVert #1\right\rVert}
\newcommand{\abs}[1]{\left\lvert #1\right\rvert}
\newcommand{\code}[1]{\nolinkurl{#1}}
\setlength{\emergencystretch}{3em}
\sloppy

\newtheorem{theorem}{Theorem}
\newtheorem{proposition}{Proposition}
\newtheorem{lemma}{Lemma}
\newtheorem{remark}{Remark}

\title{General-\(m\) Minimal-Equilibrium Stability\\
\large Proved stability regions, the bifurcation boundary, and an exact
\(m=3\) counterexample}
\author{\code{ShenWork} formalization project}
\date{27 July 2026}

\begin{document}
\maketitle

\begin{abstract}
We study the one-dimensional Neumann minimal chemotaxis system with
population mobility \(u^m\), extending the \(m=1\) stability results in
Chen--Ruau--Shen Part~II.  Three positive results are formalized: the full
two-disjunct global-stability formula for \(1\le m<2\); its critical \(m=2\)
version under an explicit endpoint admissibility condition; and a
small-sensitivity theorem for every \(m>1\), obtained from good times for the
negative-power energy and oscillation control of \(u^{1-m}\).

The energy-method boundary \(m=2\) is not a boundary of the dynamics.  A
weakly nonlinear calculation gives the direction coefficient of the first
Neumann steady branch.  Its sign changes at
\(m_c(U)\approx2.867,2.937,3.006\) for masses
\(U=0.5,1,2\), respectively.  A negative coefficient predicts a backward
branch and hence a steady obstruction below the linear threshold.  For the
fixed tuple
\[
 m=3,\qquad \beta=\gamma=\mu=\nu=U=1,
\]
this obstruction is not merely numerical: an amplitude-factored implicit
branch in a Wiener algebra, its negative direction coefficient, the positive
classical steady state, and the failure of global attraction are all
formalized in Lean 4.30.0 / Mathlib v4.30.0.

We distinguish rigorously between the proved \(m=3\) counterexample, the
formalized sufficient stability regions, and the analytically/numerically
located direction-change curve.  In particular, the sign \(m<m_c(U)\) by
itself is not claimed to prove global stability up to the linear threshold.
\end{abstract}

\tableofcontents

\section{Model, mass constraint, and the stability question}

On the unit interval \(I=(0,1)\), consider
\begin{align}
 u_t
  &=u_{xx}-\chi_0\partial_x
     \left(\frac{u^m}{(1+v)^\beta}v_x\right),
                                                         \label{eq:u}\\
 0&=v_{xx}-\mu v+\nu u^\gamma,                           \label{eq:v}\\
 u_x(t,0)&=u_x(t,1)=v_x(t,0)=v_x(t,1)=0.                 \label{eq:bc}
\end{align}
This is the minimal model \(a=b=0\).  Integrating \eqref{eq:u} and using
\eqref{eq:bc} shows that
\[
  \int_0^1u(t,x)\,dx=U
\]
is conserved.  On this mass hyperplane, the corresponding constant
equilibrium is
\begin{equation}
  u_*=U,\qquad v_*=\frac{\nu}{\mu}U^\gamma.              \label{eq:eq}
\end{equation}

The global-stability property formalized in the repository quantifies over
positive global bounded classical solutions with this conserved mass.  It
has two conclusions:
\begin{enumerate}[leftmargin=2em]
\item uniform convergence \(u(t)\to U\) in the supremum norm;
\item for each orbit, constants \(C,\rho,t_0>0\) such that the population
  and signal converge exponentially in the stated \(C^1\)-type norms for
  \(t\ge t_0\).
\end{enumerate}
The positive entry time is essential because the paper's data are not
assumed \(C^1\) at \(t=0\).

\section{Linear threshold}

Let \(\lambda_n=n^2\pi^2\), \(n\ge0\), be the Neumann eigenvalues.  Linearize
at \eqref{eq:eq}, writing \(u=U+p\), \(v=v_*+q\).  For mode \(n\ge1\),
\eqref{eq:v} gives
\[
  \widehat q_n
   =\frac{\nu\gamma U^{\gamma-1}}{\mu+\lambda_n}
      \widehat p_n.
\]
The population growth rate is therefore
\begin{equation}
 \sigma_n(\chi_0)
 =-\lambda_n+
 \chi_0
 \frac{\nu\gamma U^{m+\gamma-1}\lambda_n}
 {(1+v_*)^\beta(\mu+\lambda_n)}.                        \label{eq:sigma}
\end{equation}
For the unit interval, the threshold for the first nonzero mode is
\begin{equation}
 \chi_{\mathrm{lin}}(U,m)
 =\frac{(1+v_*)^\beta(\mu+\pi^2)}
        {\nu\gamma U^{m+\gamma-1}}.                     \label{eq:chilin}
\end{equation}
Since \(\mu+\lambda_n\) increases with \(n\), the first mode is the first to
lose stability.  Thus \(\chi_0<\chi_{\mathrm{lin}}\) is the relevant
linearly stable half-line for the unit interval.

For the tuple used in the exact counterexample,
\[
 U=\beta=\gamma=\mu=\nu=1,
\]
so \(v_*=1\) and
\begin{equation}
  \chi_{\mathrm{lin}}=2(1+\pi^2)\approx21.73920880.
  \label{eq:fixed-chilin}
\end{equation}

\section{Formalized positive stability results}

\subsection{The exact status table}

\begin{center}
\small
\renewcommand{\arraystretch}{1.2}
\begin{tabular}{@{}p{0.22\textwidth}p{0.31\textwidth}p{0.38\textwidth}@{}}
\toprule
Exponent range & Additional condition & Formalized conclusion\\
\midrule
\(1\le m<2\) &
Either disjunct of the explicit full-\(m\) global-stability formula &
Mass-constrained global sup convergence and orbitwise eventual exponential
\(C^1\) convergence\\
\(m=2\) &
The same formula plus
\(\chi_0\sqrt\mu\,\sigma_\beta U\max\{2,\gamma\}<1\) &
The critical endpoint version of the same conclusion\\
Every \(m>1\) &
Linear stability and the explicit small-sensitivity inequality associated
with a certified local basin radius &
The same global/eventual conclusion\\
\(m=3\), fixed tuple, some
\(0<\chi_0<\chi_{\mathrm{lin}}\) &
None &
Global attraction is false: an exact nonconstant positive steady orbit
exists\\
\bottomrule
\end{tabular}
\end{center}

\noindent The first two rows are sufficient conditions, not the assertion
that every linearly stable parameter is globally stable.  The fourth row is
an exact false-side theorem.  Together they show that \(m=2\) is an analytic
method boundary, not a universal dynamical threshold.

\subsection{\(1\le m<2\): entropy and signal-energy disjuncts}

The formal predicate
\code{MinimalGlobalStabilityFormula}\allowbreak\code{ConditionFullM}
packages two explicit conditions:
\[
 \text{minimal1 entropy condition}
 \quad\lor\quad
 \text{minimal2 signal-energy condition}.
\]
The minimal2 branch additionally has \(\gamma=1\).  Both thresholds include
the relevant linear-stability and nonlinear basin-entry inequalities,
evaluated using the canonical eventual upper and signal-lower constants.

\begin{theorem}[subcritical full formula]
Assume \(1\le m<2\), \(a=b=0\), \(\beta\ge1\), \(U>0\), and either
full-\(m\) formula disjunct.  Then \((U,\nu U^\gamma/\mu)\) is globally
attracting among positive bounded mass-\(U\) orbits, with orbitwise eventual
exponential \(C^1\) convergence.
\end{theorem}

\noindent This is
\code{intervalDomainM_Theorem_2_5_EventualGlobalStabilityFormula} in
\code{Paper3/IntervalDomainMMinimalSignalEnergyGlobal.lean}.  Its two
components are
\code{intervalDomainM_Theorem_2_5_minimal1_EventualGlobalStabilityFormula}
and
\code{intervalDomainM_Theorem_2_5_minimal2_EventualGlobalStabilityFormula}.

\subsection{The critical endpoint \(m=2\)}

At \(m=2\), the mass-seeded \(L^P\) recurrence becomes critical.  The missing
room is supplied by the explicit admissibility condition
\begin{equation}
 \chi_0\sqrt\mu\,\sigma_\beta\,U\max\{2,\gamma\}<1,
 \qquad
 \sigma_\beta=
 \begin{cases}
  1,&\beta=1,\\[2mm]
  \dfrac{(\beta-1)^{\beta-1}}{\beta^\beta},&\beta>1.
 \end{cases}                                           \label{eq:critical-adm}
\end{equation}
Here \(\sigma_\beta\) is the sharp upper factor for
\(z/(1+z)^\beta\) on \(z\ge0\).  Condition
\eqref{eq:critical-adm} leaves room for a restarted exponent strictly above
two and at least \(\gamma\); the endpoint Agmon estimate then supplies the
eventual box used by both formula branches.

\begin{theorem}[critical full formula]
Assume \(m=2\), \(a=b=0\), \(\beta\ge1\), \(U>0\),
\eqref{eq:critical-adm}, and either full-\(m\) formula disjunct.  Then the
same global and eventual exponential conclusion holds.
\end{theorem}

\noindent The Lean capstone is
\code{intervalDomainM_Theorem_2_5_critical_EventualGlobalStabilityFormula}
in \code{Paper3/IntervalDomainMMinimalCriticalGlobal.lean}.

\subsection{Every \(m>1\): the small-sensitivity route}

This route is independent of the \(m<2\) interpolation budget.  Its key
quantity is the negative-power weighted gradient
\begin{equation}
 G(t)=\int_0^1u(t,x)^{-2m}u_x(t,x)^2\,dx.               \label{eq:G}
\end{equation}
The entropy production provides arbitrarily late good times at which \(G\)
is as close as desired to the explicit floor
\begin{equation}
  \chi_0^2\mu\sigma_\beta^2.                            \label{eq:floor}
\end{equation}

Since
\[
  \partial_x u^{1-m}=(1-m)u^{-m}u_x,
\]
Cauchy--Schwarz on the unit interval gives the exact oscillation estimate
\begin{equation}
 \abs{u(t,x)^{1-m}-u(t,z)^{1-m}}
 \le(m-1)\sqrt{G(t)}.                                  \label{eq:osc}
\end{equation}
The mass constraint and continuity provide a point \(z\) with \(u(t,z)=U\).
Thus a good time controls \(u^{1-m}\) uniformly around \(U^{1-m}\).

For a desired local supremum radius \(0<\delta<U\), define
\begin{equation}
 r_m(U,\delta)=
 \min\left\{
  U^{1-m}-(U+\delta)^{1-m},\
  \frac{(U-\delta)^{1-m}-U^{1-m}}2
 \right\}.                                             \label{eq:r}
\end{equation}
Because \(1-m<0\), both entries are positive.  If
\begin{equation}
  (m-1)\chi_0\sqrt\mu\,\sigma_\beta<r_m(U,\delta),
  \label{eq:small}
\end{equation}
then a sufficiently late good slice satisfies
\[
  \norm{u(t,\cdot)-U}_\infty<\delta.
\]
Linear stability supplies a positive local basin radius \(\delta\); after
the good slice enters that basin, an autonomous time shift and the strong
sectorial bootstrap give eventual exponential convergence.

Equivalently, the explicit sensitivity threshold is
\begin{equation}
 \chi_0<
 \frac{r_m(U,\delta)}
 {(m-1)\sqrt\mu\,\sigma_\beta}.
 \label{eq:small-threshold}
\end{equation}

\begin{theorem}[all-\(m>1\) small-sensitivity branch]
Assume \(m>1\), \(a=b=0\), \(\beta\ge1\), \(U>0\), linear stability, and
\eqref{eq:small-threshold} for a certified local basin radius
\(0<\delta<U\).  Then the constant equilibrium is globally attracting on the
mass-\(U\) hyperplane, and every positive bounded orbit converges
exponentially in \(C^1\) after an orbit-dependent positive time.
\end{theorem}

\noindent The principal Lean statements are
\code{intervalDomainM_negativePower_oscillation_le},
\code{intervalDomainM_minimal_exists_late_supClose_smallSensitivity}, and
\code{intervalDomainM_Theorem_2_5_supercritical_smallSensitivity}.  The
theorem
\code{intervalDomainM_supercritical_smallSensitivity_exists_threshold}
constructs an admissible \(\delta\) from linear stability and exposes
\eqref{eq:small-threshold} as the remaining sufficient inequality.

\section{The steady-state bifurcation mechanism}

\subsection{Integrated flux equation}

To locate a false side, specialize first to
\[
  \beta=\gamma=\mu=\nu=1.
\]
At a positive steady state, the first equation and Neumann conditions give
\[
  u_x-\chi\frac{u^m}{1+v}v_x=0.
\]
Let \(H_m'(u)=u^{-m}\):
\[
 H_m(u)=
 \begin{cases}
  \log u,&m=1,\\[1mm]
  \dfrac{u^{1-m}}{1-m},&m\ne1.
 \end{cases}
\]
Then every positive steady state satisfies
\begin{equation}
  H_m(u)-\chi\log(1+v)=C,\qquad
  -v_{xx}+v=u,                                         \label{eq:integrated}
\end{equation}
together with \(\int_0^1u=U\).  This removes one derivative and is the
equation used by both the weakly nonlinear calculation and the finite
difference continuation.

\subsection{Weakly nonlinear expansion}

Let
\[
 e_1(x)=\cos(\pi x),\qquad e_2(x)=\cos(2\pi x),
\]
and define the elliptic response factors
\[
 r_1=\frac1{1+\pi^2},\qquad
 r_2=\frac1{1+4\pi^2},\qquad b=1+U.
\]
Seek a mass-preserving branch
\begin{align}
 u&=U+\varepsilon e_1+\varepsilon^2a_2e_2+O(\varepsilon^3), \label{eq:uexp}\\
 v&=U+\varepsilon r_1e_1+\varepsilon^2a_2r_2e_2
       +O(\varepsilon^3),                               \label{eq:vexp}\\
 \chi&=\chi_{\mathrm{lin}}+\varepsilon^2\chi_2
       +O(\varepsilon^3).                               \label{eq:chiexp}
\end{align}
The first-order solvability condition gives
\begin{equation}
 \chi_{\mathrm{lin}}=\frac{b}{r_1U^m}.                  \label{eq:weak-lin}
\end{equation}
The second harmonic at order \(\varepsilon^2\) gives
\begin{equation}
 a_2=
 \frac{m/U-r_1/b}{4(1-r_2/r_1)}.                       \label{eq:a2}
\end{equation}
Projecting the third-order equation onto \(e_1\) yields
\begin{align}
 \chi_2&=\chi_{\mathrm{lin}}S(m,U),                     \label{eq:chi2}\\
 S(m,U)
 &=\frac{a_2}{2}\left(\frac{r_2}{b}-\frac mU\right)
   +\frac{m(m+1)}{8U^2}
   -\frac{r_1^2}{4b^2}.                                \label{eq:S}
\end{align}
At the level of this weakly nonlinear expansion:
\[
  S>0\ \Longrightarrow\ \text{the local branch points to }
       \chi>\chi_{\mathrm{lin}},
\]
whereas
\[
  S<0\ \Longrightarrow\ \text{the branch is backward and contains }
       \chi<\chi_{\mathrm{lin}}.
\]
If the formal branch is promoted to an actual nonconstant positive steady
branch, every point on its backward side is a bounded global orbit that
cannot converge to the constant equilibrium.  The required analytic branch
construction is supplied below for the fixed \(m=3,U=1\) tuple; no
all-\(m\) branch theorem is being asserted here.

\subsection{Direction-change curve}

Solving \(S(m,U)=0\) gives:
\begin{center}
\renewcommand{\arraystretch}{1.15}
\begin{tabular}{@{}ccc@{}}
\toprule
Mass \(U\) & \(m_c(U)\) & Sign for \(m>m_c(U)\)\\
\midrule
\(0.5\) & \(2.866754572143\) & backward\\
\(1\)   & \(2.936658202934\) & backward\\
\(2\)   & \(3.005668668328\) & backward\\
\bottomrule
\end{tabular}
\end{center}
These numbers are roots of the explicit coefficient
\eqref{eq:S}, not fitted continuation parameters.

\begin{remark}[scope of the curve]
The curve \(S=0\) is a local steady-bifurcation direction boundary.  For
\(m>m_c(U)\), the coefficient predicts a direct obstruction below the linear
threshold.  The required local branch theorem and the resulting obstruction
are carried out completely in Lean for the fixed \(m=3,U=1\) tuple.  For
\(m<m_c(U)\), a positive coefficient merely removes this predicted local
obstruction.  It does not prove global attraction for every
\(0<\chi_0<\chi_{\mathrm{lin}}\).
\end{remark}

\section{Numerical cross-check}

The script \code{research_notes/m_supercritical_bifurcation_probe.py}
performs two independent checks.

\paragraph{Discrete weak coefficient.}
For \(N\) finite-difference subintervals it replaces
\[
 \pi^2,\ 4\pi^2
\quad\text{by}\quad
 \lambda_{1,h}=\frac4{h^2}\sin^2\frac{\pi}{2N},\qquad
 \lambda_{2,h}=\frac4{h^2}\sin^2\frac{\pi}{N}
\]
in \eqref{eq:a2}--\eqref{eq:S}.

\paragraph{Nonlinear solve.}
At every grid node it enforces the integrated flux equation
\eqref{eq:integrated}, the finite-difference elliptic equation with reflected
Neumann ghost points, the trapezoidal mass constraint, and a prescribed first
cosine coefficient.  It then compares the fitted quadratic branch direction
with the discrete analytic coefficient.

For \(N=160\):
\begin{center}
\small
\renewcommand{\arraystretch}{1.15}
\begin{tabular}{@{}cccccc@{}}
\toprule
\(m\)&\(U\)&\(\chi_{\mathrm{lin},h}\)&analytic \(\chi_{2,h}\)&fitted
\(\chi_{2,h}\)&largest residual\\
\midrule
4&1&21.7385746366&\(-4.24003585072\)&\(-4.24003584422\)&\(1.1\times10^{-11}\)\\
3&0.5&130.431447819&\(-9.59085815306\)&\(-9.59085814853\)&\(4.8\times10^{-12}\)\\
\bottomrule
\end{tabular}
\end{center}
At relative amplitude \(0.02\), the sensitivity shifts are
\[
 \Delta\chi=-1.6966\times10^{-3}\quad(m=4,U=1),
\]
and
\[
 \Delta\chi=-9.5917\times10^{-4}\quad(m=3,U=0.5).
\]
Pseudo-arclength continuation preserves the backward direction through the
reported steps.  These computations validate the coefficient and the
implementation, but numerical residuals are not used as proof of the exact
\(m=3\) counterexample.

\section{Exact Lean proof of the \(m=3\) false side}

\subsection{Functional setting}

Fix
\[
  m=3,\qquad \beta=\gamma=\mu=\nu=U=1.
\]
The steady system is
\begin{align}
 0&=u_{xx}-\chi\partial_x\left(\frac{u^3v_x}{1+v}\right),
                                                               \label{eq:m3u}\\
 0&=v_{xx}-v+u.                                                \label{eq:m3v}
\end{align}
The formalization works in a weighted \(\ell^1\) Wiener algebra of Fourier
coefficients.  This choice has four useful properties:
\begin{enumerate}[leftmargin=2em]
\item multiplication is continuous, so nonlinear products remain in the
  algebra;
\item the coefficient norm controls pointwise evaluation;
\item two Wiener weights control the two spatial derivatives needed for a
  classical steady state;
\item even real sequences encode Neumann cosine profiles, and the mean-zero
  subspace enforces the mass constraint.
\end{enumerate}

Split the mean-zero even profile into its first cosine mode and a closed
complement:
\[
  h=e_1+z,\qquad z\perp e_1.
\]
With amplitude \(a\),
\[
  u_a=1+ah,\qquad v_a=1+aRh,
\]
where \(R=(-\partial_{xx}+1)^{-1}\) is the diagonal Fourier resolver.
Dividing the steady residual by the amplitude produces a smooth blown-up
map in \((a,z,\chi)\).  At
\[
  (a,z,\chi)=(0,0,2(1+\pi^2))
\]
its derivative in \((z,\chi)\) is a bounded linear isomorphism.  The Banach
implicit function theorem therefore yields a \(C^3\) branch
\[
  a\longmapsto(z(a),\chi(a)).
\]
The construction and inverse are in
\code{MinimalSteadyWABlowup.lean},
\code{MinimalSteadyWALinearIso.lean}, and
\code{MinimalSteadyWAImplicitBranch.lean}.

\subsection{Exact branch direction}

After multiplying out the harmless signal denominator, the branch satisfies
the polynomial flux identity
\begin{equation}
  (2+aRh)\,Dh=\chi(a)(1+ah)^3DRh.                       \label{eq:poly}
\end{equation}
The mode-two component of the first derivative determines the second-harmonic
correction.  The mode-one component of the second derivative gives
\begin{equation}
 \chi'(0)=0,\qquad
 \chi''(0)=
 -\frac{6\pi^4+37\pi^2+25}
 {24\pi^2(\pi^2+1)}<0.                                 \label{eq:exact-second}
\end{equation}
Since the expansion convention \eqref{eq:chiexp} has
\(\chi_2=\chi''(0)/2\), \eqref{eq:exact-second} equals
\[
  \chi_2\approx-0.189271408,
\]
exactly matching \eqref{eq:S} at \(m=3,U=1\).

The formal theorems are
\code{minimalSensitivityFirstJet_eq_zero},
\code{minimalSensitivitySecondJet_eq}, and
\code{minimalSensitivitySecondJet_neg} in
\code{Paper3/MinimalSteadyWADirection.lean}.  Elementary derivative tests
then show that \(\chi(a)\) is strictly decreasing for all sufficiently small
\(a>0\), while continuity keeps it positive:
\begin{equation}
  0<\chi(a)<2(1+\pi^2).                                 \label{eq:backward}
\end{equation}
This selection is formalized in
\code{Paper3/MinimalSteadyWABranchSelection.lean}.

\subsection{From the Banach branch to a classical counterexample}

For sufficiently small positive \(a\), the Wiener norm bound gives
\[
  u_a(x)>0,\qquad v_a(x)>0\quad\text{for all }x.
\]
Coefficient evaluation and Fourier differentiation turn the algebraic
residual into the pointwise classical equations
\eqref{eq:m3u}--\eqref{eq:m3v}.  Evenness gives the Neumann boundary
conditions, the missing zero Fourier coefficient gives
\[
  \int_0^1u_a(x)\,dx=1,
\]
and the fixed first coefficient makes \(u_a\) nonconstant.

\begin{theorem}[exact subcritical steady state]
There exist \(0<\chi<2(1+\pi^2)\) and positive, nonconstant, \(C^2\) Neumann
profiles \(u,v\) satisfying \eqref{eq:m3u}--\eqref{eq:m3v} and
\(\int_0^1u=1\).
\end{theorem}

\noindent This is
\code{exists_nonconstant_minimal_steady_below_threshold_m3} in
\code{Paper3/MinimalSteadyWACounterexample.lean}.

View the profiles as constant-in-time functions.  They form a positive,
bounded, mass-one global classical orbit.  If the constant equilibrium
\((1,1)\) attracted every such orbit, then \(u(x)\) would converge both to
its constant-in-time value \(u(x)\) and to \(1\).  Uniqueness of limits would
force \(u\equiv1\), contradicting nonconstancy.

\begin{theorem}[global attraction is false below the linear threshold]
For the fixed \(m=3\) tuple, there exists
\(0<\chi<\chi_{\mathrm{lin}}\) such that the minimal equilibrium does not
attract all positive bounded mass-one classical orbits.
\end{theorem}

\noindent The exact Lean theorem is
\code{minimal_equilibrium_global_stability_false_m3}.  This completes the
false side inside the proof assistant; no unformalized
Crandall--Rabinowitz theorem or numerical existence assumption remains in
this tuple.

\section{What is resolved, and what is not claimed}

The results establish the following.
\begin{enumerate}[leftmargin=2em]
\item The \(m<2\) energy method extends the paper's \(m=1\) result to
  \(1\le m<2\) under the full explicit formula conditions.
\item The critical exponent \(m=2\) is included under the explicit endpoint
  admissibility \eqref{eq:critical-adm}.
\item For every \(m>1\), sufficiently small sensitivity in the explicit
  sense \eqref{eq:small-threshold} gives global attraction, provided the
  equilibrium is linearly stable.
\item The assertion ``every linearly stable minimal equilibrium is globally
  attracting for every \(m>1\)'' is false.  It already fails at \(m=3\) for
  a sensitivity strictly below \(\chi_{\mathrm{lin}}\).
\item The direction curve \(m_c(U)\) explains why the false side begins near
  \(m\approx2.9\) in the normalized family and why the energy boundary
  \(m=2\) is not sharp dynamically.
\end{enumerate}

\noindent What is \emph{not} proved is a complete necessary-and-sufficient
global-stability diagram in the region between the sufficient thresholds and
the subcritical steady branch.  In particular:
\begin{itemize}[leftmargin=2em]
\item \(m<m_c(U)\) does not imply global stability all the way to
  \(\chi_{\mathrm{lin}}\);
\item \(m>m_c(U)\) makes the weak coefficient negative; the corresponding
  local steady obstruction is proved here only for the fixed \(m=3,U=1\)
  tuple and does not classify every smaller sensitivity;
\item the three numerical values of \(m_c(U)\) belong to the normalized
  \(\beta=\gamma=\mu=\nu=1\) family, not to arbitrary parameter tuples.
\end{itemize}
These qualifications are part of the result: they separate exact theorems
from the phase diagram suggested by bifurcation analysis.

\section{Lean theorem map}

\begin{center}
\small
\renewcommand{\arraystretch}{1.18}
\begin{tabular}{@{}p{0.34\textwidth}p{0.57\textwidth}@{}}
\toprule
Component & Principal source / theorem\\
\midrule
\(1\le m<2\), full formula &
\code{IntervalDomainMMinimalSignalEnergyGlobal.lean}:
\code{intervalDomainM_Theorem_2_5_EventualGlobalStabilityFormula}\\
\(m=2\), critical endpoint &
\code{IntervalDomainMMinimalCriticalGlobal.lean}:
\code{intervalDomainM_Theorem_2_5_critical_EventualGlobalStabilityFormula}\\
All \(m>1\), negative-power good slice &
\code{IntervalDomainMNegativePowerBox.lean}:
\code{intervalDomainM_minimal_exists_late_supClose_smallSensitivity}\\
All \(m>1\), small-sensitivity capstone &
\code{IntervalDomainMSupercriticalSmallSensitivity.lean}:
\code{intervalDomainM_Theorem_2_5_supercritical_smallSensitivity}\\
\(m=3\) Banach branch and direction &
\code{MinimalSteadyWAImplicitBranch.lean},
\code{MinimalSteadyWADirection.lean},
\code{MinimalSteadyWABranchSelection.lean}\\
Exact classical steady state &
\code{MinimalSteadyWACounterexample.lean}:
\code{exists_nonconstant_minimal_steady_below_threshold_m3}\\
Failure of global attraction &
\code{MinimalSteadyWACounterexample.lean}:
\code{minimal_equilibrium_global_stability_false_m3}\\
Analytic/numerical cross-check &
\code{research_notes/m_supercritical_bifurcation_probe.py}\\
\bottomrule
\end{tabular}
\end{center}

\medskip
\noindent All Lean capstones above are \code{sorry}-free.  Their printed axiom
audits use only the standard classical axioms
\code{propext}, \code{Classical.choice}, and \code{Quot.sound}.  The numerical
script is reproducible but is supplementary; the \(m=3\) false-side theorem
does not depend on it.

\end{document}
